\documentclass[11pt]{article}
\input{preamble}
\setdoc{Method}{How we work: evidence rules, bias control, phases}{v0.1}{Draft -- EDA in progress}{2026-11-27}

\begin{document}
\makedoctitle
\tableofcontents
\clearpage

% =====================================================================
\section{Purpose}
This document gives the procedure that the team follows from framing to delivery.
The procedure is more important than any single result: a result is only as good as the procedure that produced it.
Each phase has questions, inputs, activities, a deliverable and an exit criterion.

% =====================================================================
\section{Evidence rules}
\subsection{Evidence tags}
Each statement of fact in the project documents has one tag.

\taglegend

\subsection{Falsification instead of confirmation}
\begin{enumerate}
  \item Write each hypothesis as a claim that a test can refute.
  \item Write a rival hypothesis for each claim. The rival is a real alternative, not a weak version of the claim.
  \item Write the refutation test and the result that refutes the claim \emph{before} you run the test.
  \item Run the test. Record the result in the register, also if the claim is refuted.
  \item Use the status \status{NOT REFUTED} when the test does not refute the claim. Never use ``confirmed''.
  \item If a test cannot fail, it is not a test. Change the test or keep the status \status{OPEN} and say why.
\end{enumerate}

\subsection{Decision lifecycle with rivals}
Each decision has the status \sWorking{} and at least one rival with the status \sProposed.

\begin{center}
\begin{tikzpicture}[node distance=4.6cm, every node/.style={font=\small\sffamily}]
  \node[draw, rounded corners, fill=projblue!8] (w) {WORKING};
  \node[draw, rounded corners, fill=tagunknown!8, right of=w] (r) {REFUTED};
  \node[draw, rounded corners, fill=projblue!8, right of=r] (n) {new WORKING};
  \node[draw, rounded corners, fill=gray!10, below of=w, yshift=3cm] (p) {PROPOSED (rival)};
  \draw[->] (w) -- node[above, font=\scriptsize] {evidence negated} (r);
  \draw[->] (r) -- node[above, font=\scriptsize] {evaluate rivals} (n);
  \draw[->, dashed] (p) -| node[below, pos=0.25, font=\scriptsize] {evaluated} (n);
\end{tikzpicture}
\end{center}

\begin{enumerate}
  \item A working decision stays in use while the evidence for it is not negated.
  \item Each decision has a written switch condition (the ``refutation'' field in the register).
  \item When the switch condition occurs, set the decision to \sRefuted. Evaluate the rivals with the same rules. Record the new decision with a new ID and a link to the old one.
  \item Do not delete a rival when it is not used. A rival can become the working decision later.
\end{enumerate}

\subsection{Fix the rules before the results}
Before modelling, fill the pre-registration sheet (Section~\ref{sec:prereg}). Sign and date it. Commit it.
After this, these changes are logged with a reason and a date: metric, baseline, split date, threshold method, failure line, hyperparameter search space.
The test period is used one time, at the end. If it is used again, the report says so and gives the reason.

% =====================================================================
\section{Bias checklist}
Use this list at the end of each phase. Write the result in the phase deliverable.

\small
\begin{longtable}{@{}P{0.18\linewidth} P{0.38\linewidth} P{0.38\linewidth}@{}}
\toprule
\textbf{Bias} & \textbf{Where it can enter this project} & \textbf{Counter-measure} \\
\midrule\endhead
\bottomrule\endfoot
Confirmation & We look only for evidence that supports the product idea. & Rival for each claim. Search for disconfirming sources (\refid{S-07}, \refid{S-14}). \\
Anchoring & The first threshold (75th percentile) or the first model becomes the default. & Rival thresholds in \refid{D-003}. Sensitivity analysis. \\
HARKing (hypotheses after results) & A pattern found in EDA becomes a ``hypothesis'' that it already supports. & Hypotheses from EDA are tested on a period that EDA did not use, or are marked as exploratory. \\
Forking paths & Many feature sets or models are tried; only the best one is reported. & Pre-registration. Log every run. Report the number of tries. \\
Look-ahead and target leakage & Rolling means include the target hour. Threshold uses test data. Imputer fit on all data. & Past-only windows. Fit every stage on training data. Leakage check (Section~\ref{sec:leak}). \\
Selection and survivorship & 4 stations chosen by the producer. Periods with outages are dropped silently. & Note in the 5W2H. \refid{G-17}, \refid{H-05}. Flags instead of deletion (\refid{D-009}). \\
Availability & We use only the data we already have because it is easy. & Source search (\refid{D-015}). \\
Sunk cost & We keep a framing because we worked on it. & Switch conditions written before the work (\refid{D-001}). \\
Narrative & The talk tells a clean story and hides weak points. & Talk shows baselines, per-station results, refuted hypotheses. \\
Automation & People trust model output over data quality. & Data-trust flag shown next to each class. \\
\end{longtable}
\normalsize

% =====================================================================
\section{Phases}
The phases follow CRISP-DM, adapted to the course brief and to the README checklist.
A later phase can open an earlier one again. When this occurs, record the reason in the decision log.

\newcommand{\phase}[6]{%
\subsection{#1}
\begin{tabularx}{\linewidth}{@{}l X@{}}
\textbf{Questions} & #2 \\
\textbf{Inputs} & #3 \\
\textbf{Activities} & #4 \\
\textbf{Deliverable} & #5 \\
\textbf{Exit criterion} & #6 \\
\end{tabularx}}

\phase{Phase 1 -- Frame}
{Who needs what, and why? What can refute the idea?}
{Course brief, first look at the data.}
{Write the charter, the 5W2H, the first gaps, hypotheses and decisions with rivals.}
{\fpath{product-charter.pdf}, registers.}
{Each working decision has a rival and a switch condition.}

\phase{Phase 2 -- Discover sources and understand the data}
{Where does the data come from? What does one value mean? What other data exists?}
{Source register, producer pages, reports.}
{Trace the data lineage. Search the portals. Search for disconfirming sources. Fill the source grid. Write the raw file manifest.}
{\fpath{data.pdf} part 1, updated gaps.}
{Gaps \refid{G-01} to \refid{G-03} and \refid{G-13} are closed, or their impact is written.}

\phase{Phase 3 -- Audit and flag}
{Which values are doubtful, and why? Is missing data random?}
{Raw files (read only), manifest.}
{Write Bronze (typed copy). Run the quality checks of the audit protocol. Add flags in Silver. Do not delete.}
{Silver data in \fpath{spark/lake/silver/}, audit results in \fpath{data.pdf} part 2.}
{Each quality rule has a count per station and a linked gap.}

\phase{Phase 4 -- Explore}
{What patterns exist? What alternative explanations exist for each one?}
{Silver data.}
{Profiles per station and period. For each pattern, write at least one alternative explanation. Use the training period only for patterns that become hypotheses.}
{Notebook sections, new or updated hypotheses.}
{Each interpretation in the notebook has an alternative and a tag.}

\phase{Phase 5 -- Fix the evaluation protocol}
{How will we decide if a model is useful?}
{Registers, \refid{D-003} to \refid{D-007}.}
{Fill the pre-registration sheet. Sign, date and commit it.}
{Pre-registration sheet.}
{Commit exists before the first model run.}

\phase{Phase 6 -- Build features}
{Which inputs are available at prediction time for each mode?}
{Silver data, pre-registration.}
{Calendar features. Lags and rolling means with past values only. Labels from the training threshold.}
{Gold data in \fpath{spark/lake/gold/}.}
{Leakage check passes (Section~\ref{sec:leak}).}

\phase{Phase 7 -- Model}
{Do the models do better than the baselines on validation?}
{Gold data.}
{Fit the baselines first. Fit the Spark ML Pipelines (both algorithms, both modes). Tune on the validation period only. Log every run.}
{Run log, chosen configuration for each model.}
{Configurations frozen and committed.}

\phase{Phase 8 -- Evaluate}
{What is the result on data that no step has used?}
{Frozen configurations, test period.}
{Run the evaluation one time. Compute metrics per station. Compare with baselines. Update the hypothesis statuses.}
{Results table, updated registers.}
{Every pre-registered test has a status.}

\phase{Phase 9 -- Communicate}
{What can a user trust, and what not?}
{Results, registers.}
{Write the talk (7 minutes): problem (1), data and gaps (1.5), pipeline (2), results with baselines and refuted hypotheses (1.5), limits and next steps (1). Prepare answers for 5 minutes of questions.}
{Slides, notebook, updated documents.}
{Go or no-go criteria of the charter are met.}

% =====================================================================
\section{Pre-registration sheet}
\label{sec:prereg}
Fill this sheet before Phase 7. Copy it into the notebook as the first markdown cell of the modelling section.

\begin{tabularx}{\linewidth}{@{}l X@{}}
\toprule
\textbf{Item} & \textbf{Value (fill before modelling)} \\
\midrule
Date and commit & \\
Hypotheses tested & \refid{H-01}, \refid{H-02}, \ldots{} with their refutation lines \\
Target and threshold method & \refid{D-003} \\
Horizons & \refid{D-004} \\
Split dates & \refid{D-005} \\
Primary metric & \refid{D-006} \\
Guardrail metrics & \refid{D-006} \\
Baselines & \refid{D-006} \\
Margins (minimum gain over baseline) & \\
Failure line & \\
Algorithms and search space & \refid{D-007} \\
Flags excluded from training & \refid{D-009} \\
Seeds & \\
\bottomrule
\end{tabularx}

% =====================================================================
\section{Leakage rules}
\label{sec:leak}
\begin{enumerate}
  \item Split by time before any fit. Training before validation before test.
  \item Compute the class threshold on training data only.
  \item Fit every Spark ML stage (\texttt{Imputer}, \texttt{StringIndexer}, \texttt{StandardScaler}, models) on training data only, inside the Pipeline.
  \item Use window frames that end at the hour before the prediction time: \texttt{rowsBetween(-k, -1)} or a range frame on time.
  \item For horizon $h$, use only inputs that exist at time $t$ to predict $t+h$.
  \item Do not fill missing values with data from the future (no backward fill, no centred interpolation).
  \item Check: shuffle the target in time. If the model still does better than the baseline, look for leakage.
\end{enumerate}

% =====================================================================
\section{Data handling rule}
\begin{enumerate}
  \item Treat \fpath{spark/data/} as a raw zone. Do not write, edit, filter or overwrite files there.
  \item Add a file to the raw zone only with a documented download step. Write a manifest entry (see \fpath{data.pdf}).
  \item Write all processed data to \fpath{spark/lake/} (Bronze, Silver, Gold).
  \item Flag doubtful values. Do not delete them.
  \item Make each processing step reproducible from Bronze.
\end{enumerate}
\textbf{Proposed enforcement} \sProposed: mount \fpath{./data} as read-only and add a writable \fpath{./lake} mount in \fpath{spark/docker-compose.yml}. This change needs team approval.

% =====================================================================
\section{Traceability}
Each item of work links a need to a decision. Use the IDs in notebook markdown cells and in commit messages.

\begin{tabularx}{\linewidth}{@{}X X X X X@{}}
\toprule
\textbf{Need} & \textbf{Question} & \textbf{Gap or hypothesis} & \textbf{Analysis} & \textbf{Decision} \\
\midrule
Warning before exposure & Are levels predictable? & \refid{H-01}, \refid{H-02} & Notebook, modelling & \refid{D-001}, \refid{D-004} \\
Meaningful class & What is ``high''? & \refid{G-14}, \refid{H-06} & Sources, Silver counts & \refid{D-003} \\
Trust in data & Which values are doubtful? & \refid{G-07} to \refid{G-11}, \refid{H-07} & Notebook section 4 & \refid{D-009} \\
Fair evaluation & Which metric fits the users? & \refid{G-15} & Sources, users & \refid{D-006} \\
Generalisation & Do results apply to other stations? & \refid{G-17}, \refid{H-05} & Leave-one-station-out & \refid{D-001} \\
Course brief & Is the chain complete? & -- & Notebook & \refid{D-007}, \refid{D-008} \\
\bottomrule
\end{tabularx}

% =====================================================================
\section{Reproducibility}
\begin{itemize}
  \item Read CSV files with an explicit schema. Do not infer types.
  \item Set the Spark session time zone to UTC.
  \item Fix all random seeds. Keep \texttt{PYTHONHASHSEED=1} (\fpath{spark/Dockerfile}).
  \item Run the full notebook from a clean clone in the container before each delivery.
  \item Record the git commit and the raw file checksum with each result.
\end{itemize}

\clearpage
\input{glossary}
\end{document}
